\section{Topics on Password}
\label{sec::chp2::password}

Passwords are obvious user secrets. Much of this dissertation focuses on human-chosen and memorable secrets, such as \password{il0vesk00l!}, as their weaknesses stem from predictable and frequently used patterns. Password security and authentication constitute a vast field encompassing multiple subfields. Some examples are as follows.


\begin{itemize}
%[itemsep=0pt, parsep=0pt, leftmargin=10pt]
    \item Password usability and user experience
        \begin{itemize}
            \item Password security in enterprise environments
            \item Password creation and memorability
            \item User education and awareness
            \item Password statistics (length, count/freq, structure info)
        \end{itemize}
    \item Password reset and recovery mechanisms
    \item Password modelling and similarity
        \begin{itemize}
            \item Similarity (e.g. distance-like metrics)
        \end{itemize}
    \item Password entropy and complexity analysis 
        \begin{itemize}
            \item Password composition policy (PCP) (or password creation policy)
            \item Password strength meters (PSM)
        \end{itemize}
    \item Password guessing algorithms
    \item Password hashing and encryption techniques
    \item Password-less authentication schemes
    \item Password authentication protocols
    \item Federated identity and single sign-on (SSO)
    \item Biometric and multi-factor authentication
    \item Password-based authentication in specific domains (e.g., web, mobile, IoT)
    \item Password managers and storage
\end{itemize}


\textbf{Statistical and Lexical Password Analysis.} Historical analyses of password characteristics have focused on \textit{length}, \textit{count}, and \textit{character-level structures}~\cite{ma_study_2014, Li_Han_Xu_2014, Das_Bonneau_Caesar_Borisov_Wang_2014, ur_et_al_2015, li_study_2016, pearman_et_al_2017, Pasquini_Cianfriglia_Ateniese_Bernaschi_2021, Cazier_Medlin_2006}. Our research advances this field by examining nuanced aspects of passwords with greater precision. Moreover, the evolution of language models, combined with our successful prototypes, presents a cost-effective approach for labelling additional passwords.

In the analysis of Chinese passwords~\cite{Li_Han_Xu_2014}, the authors conduct a broad examination of password syntax. Our research extends~\cite{Li_Han_Xu_2014} by delving deeper into password structures. For instance, they identified $LLLLLL$ as the predominant structure within the \texttt{Rockyou} dataset, typically representing a word, as demonstrated in our results. Further structural analysis across various password datasets is presented in~\cite{li_study_2016}. Das~\etal~\cite{Das_Bonneau_Caesar_Borisov_Wang_2014} explored password syntax, notably introducing an analysis of word or phrase transformations, which aligns with aspects of our research. However, we disagree with their classification of sequential keys, alternate keys, and sequential alphabets as transformations, instead viewing these as sequences. Riddle~\etal~\cite{Riddle_Miron_Semo_1989} investigate the composition and semantics of passwords, with a particular emphasis on the psychological significance of words.

In ``Investigating the Distribution of Password Choices''~\cite{Malone_Maher_2012}, the analysis of the \hmail dataset concluded that its distribution closely resembles Zipf's law. Our analytical efforts extend this observation by demonstrating that extracted words from unique passwords also follow a similar distribution---more of this is found in Chapter~\ref{chp4::password-composition}.

\textbf{Password Modelling.} Early password modelling techniques, such as the Markov $n$-gram model~\cite{narayanan_fast_2005}, focused on individual characters, positing that \textit{``the distribution of letters in easy-to-remember passwords likely mirrors the letter distribution in the user's native language.''} We validate this in Section~\ref{chp4::password-composition} at the word level. The observation that letter distribution is not uniform holds true for words, exhibiting a power-law distribution akin to Zipf's law. Modelling passwords based on Probabilistic Context-Free Grammars (PCFG)~\cite{weir_password_2009} was the next innovation. Our insights into empirical password structures (\S\ref{sec::chp4::results}) can refine \pcfg modelling. Incorporating $W_n \text{ and } N_n$ as variables for words and names, respectively, enhances the computational viability of modelling more intricate structures. Subsequent advancements have leveraged machine learning techniques in password modelling, including neural networks~\cite{Castro_Lang_Liu_Mata_2017}, GANs~\cite{hitaj_passgan_2019}, RNNs~\cite{Melicher_Ur_Segreti_Komanduri_Bauer_Christin_Cranor_2016}, and Transformers~\cite{Xu_Yu_Zhang_Wang_Zhang_Wu_Han_2023}, among others.

\textbf{Password Composition Policy (PCP).} Despite extensive research into password complexity measures and meters, recommendations for password policies remain scarce. Bonk~\etal~\cite{Bonk_Parish_et_al_2021} provide guidance on constructing passphrases, effectively extending the NCSC's~\cite{ncsc2024} three-word suggestion to longer sequences, such as seven words or more. \zxcvbn employs dictionaries and bespoke rules to gauge password strength on a scale from 0 to 4. However, \zxcvbn has its limitations; for instance, it incorrectly parses the password \password{gomythsun} as \texttt{"go myths un"}, which affects its strength rating. Additionally, it struggles with transformations such as \texttt{"numbr"} from \texttt{"number"}. Wang~\etal~\cite{Wang_Shan_Dong_Shen_Jia_2023} conduct a deeper evaluation of password strength meters. Their endorsement of \zxcvbn and Markov $n$-gram models has informed our research approach.

The deployment of language models, particularly Large Language Models (LLMs), in this domain remains notably scarce. Rando~\etal~\cite{Rando_Perez-Cruz_Hitaj_2023} explored password modelling using the GPT-2 framework. Given the broad embeddings that encompass multiple languages, the semantic breakdown and analysis of passwords represent, in our view, a pioneering application of LLMs.
